\documentclass[11pt]{article}
\usepackage[margin=1.1in]{geometry}
\usepackage{amsmath,amssymb,amsthm}
\usepackage{enumitem}
\usepackage[colorlinks=true,linkcolor=blue,citecolor=blue,urlcolor=blue]{hyperref}

\newtheorem{theorem}{Theorem}[section]
\newtheorem{lemma}[theorem]{Lemma}
\newtheorem{proposition}[theorem]{Proposition}
\newtheorem{corollary}[theorem]{Corollary}
\newtheorem{remark}[theorem]{Remark}
\theoremstyle{definition}
\newtheorem{definition}[theorem]{Definition}

\newcommand{\pP}{\textbf{[P]}}
\newcommand{\pC}{\textbf{[C]}}
\newcommand{\pO}{\textbf{[O]}}
\newcommand{\Tr}{\operatorname{Tr}}

\title{\textbf{Condition NL for Finite Balls:\\
A Spectral Reduction, the Unconditional Metric Sector,\\
and a Universal Kinematic Residue}}
\author{Robert W.\ McGwier\\
\small CTO and Staff Scientist, Cohere Technology Group, Herndon, VA}
\date{August 7, 2026 --- Version 2}

\begin{document}
\maketitle

\begin{abstract}
\noindent
The companion note \cite{McGwier2026b} reduced the finite-ball torsionful
first law to a single condition (Condition NL): vanishing of the cross
expectation of the nonlocal modular remainder against state perturbations
dual to linearized Einstein--Cartan solutions. We attack that condition.
First, we derive from the exact operator identity
$\delta K=-\delta\log\rho$ a closed-form modular smearing representation,
$\delta K=-\int ds\,\tfrac{\pi}{2\cosh^2(\pi s)}\,
\sigma_s(\rho^{-1/2}\delta\rho\,\rho^{-1/2})$, with spectral multiplier
$\kappa/(2\sinh(\kappa/2))$, and machine-verify it to $10^{-14}$ by four
independent computations. Second, using this spectral form we prove the
principal new result: by the conformal selection rule
$\langle J_5\,T\rangle=0$, the graviton channel of Condition NL vanishes
identically, so \emph{the metric sector of the torsionful first law holds
unconditionally for finite balls}---strengthening Theorem 5.2 of
\cite{McGwier2026b} from conditional to unconditional in that sector.
Third, we show the sole surviving channel (axial--axial) is proportional
to the current central charge $C_J$ times a purely kinematic integral,
universal across all CFTs with a conserved axial current; its value is the
single number on which the remaining condition depends. If it vanishes,
Condition NL holds outright; if not, we prove it is absorbable as a finite
renormalization of the local term---an improvement of the holographic
dictionary---so the torsionful first law survives in improved form either
way. The kinematic integral itself is stated explicitly and left as the
one open computation. Labels: \pP{} proved/derived, \pC{} conjectured with
support, \pO{} open.
\end{abstract}

\section{The Problem}

Setting and notation follow \cite{McGwier2026a,McGwier2026b}: holographic
CFT$_4$, massless minimally coupled Dirac matter, theory deformed by a
smooth compactly supported spin-current source, which by the Reduction
Lemma of \cite{McGwier2026b} is a background axial field,
$\delta S_{\mathrm{def}}=\int\beta_\mu J_5^\mu$. Working at joint order
$O(\beta\lambda)$ ($\beta$ the source, $\lambda$ the state perturbation),
the deformed ball modular Hamiltonian splits as
$K_B^{(\beta)}=K_B+\delta K_B^{\mathrm{loc}}+\delta K_B^{\mathrm{nl}}$,
and:

\begin{definition}[Condition NL, restated from \cite{McGwier2026b}]
For all balls $B$ and all state perturbations $|\delta\Psi\rangle$ dual to
linearized Einstein--Cartan solutions,
$\langle0|\delta K_B^{\mathrm{nl}}|\delta\Psi\rangle+\mathrm{c.c.}=0$
at order $O(\beta\lambda)$.
\end{definition}

Under Condition NL the finite-ball torsionful first law is equivalent to
linearized EC gravity with torsion boundary data; on null planes the
condition was proved to hold via the Markov property
\cite{CasiniTesteTorroba2017}. The finite-ball case is the subject of this
paper.

\section{The Kernel, Derived and Machine-Verified}

\begin{proposition}[Exact perturbation of the modular Hamiltonian]
\label{prop:kernel}
Let $\rho=\rho_0+\delta\rho$ with $\rho_0>0$, and $K=-\log\rho$. Then at
first order, in the eigenbasis $\rho_0|i\rangle=\lambda_i|i\rangle$,
\begin{equation}
\langle i|\delta K|j\rangle
=-\,\delta\rho_{ij}\,
\frac{\log\lambda_i-\log\lambda_j}{\lambda_i-\lambda_j}
=-\int_0^\infty\! db\;
\big[(\rho_0+b)^{-1}\delta\rho\,(\rho_0+b)^{-1}\big]_{ij}\,,
\label{eq:DK}
\end{equation}
and equivalently, with modular flow
$\sigma_s(X)=\rho_0^{\,is}X\rho_0^{-is}$ and
$X=\rho_0^{-1/2}\,\delta\rho\,\rho_0^{-1/2}$,
\begin{equation}
\boxed{\;\delta K
=-\int_{-\infty}^{\infty}\! ds\;
\frac{\pi}{2\cosh^2(\pi s)}\;\sigma_s\!\big(X\big)\;}
\qquad\text{with}\qquad
\int_{-\infty}^{\infty}\!ds\,
\frac{\pi\,e^{i\kappa s}}{2\cosh^2(\pi s)}
=\frac{\kappa}{2\sinh(\kappa/2)}\,,
\label{eq:kernel}
\end{equation}
where $\kappa_{ij}=\log(\lambda_i/\lambda_j)$ is the modular frequency.
\end{proposition}

\begin{proof}
The first equality in \eqref{eq:DK} is the Daleckii--Krein derivative of
$-\log$; the second follows from
$\int_0^\infty db\,[(\lambda_i+b)(\lambda_j+b)]^{-1}
=(\log\lambda_i-\log\lambda_j)/(\lambda_i-\lambda_j)$. For
\eqref{eq:kernel}: $\sigma_s(X)_{ij}=X_{ij}e^{is\kappa_{ij}}$ and
$X_{ij}=\delta\rho_{ij}\lambda_i^{-1/2}\lambda_j^{-1/2}$, so matching to
\eqref{eq:DK} requires a kernel with Fourier transform
$\kappa\lambda_i^{1/2}\lambda_j^{1/2}/(\lambda_i-\lambda_j)
=\kappa/(2\sinh(\kappa/2))$; the stated $\cosh^{-2}$ profile has exactly
this transform. \pP{}
\end{proof}

\begin{remark}[Machine verification]
\label{rem:verification}
All identities in Proposition~\ref{prop:kernel} were verified numerically
on random $5\times5$ density matrices by four independent computations of
$\delta K$: symmetric finite difference of $-\log\rho$; the
Daleckii--Krein formula; the $b$-integral; and the modular smearing
\eqref{eq:kernel}. Maximal deviations: $1.7\times10^{-8}$ against finite
differences (consistent with $O(\epsilon^2)$ step error at
$\epsilon=10^{-6}$) and $1.3\times10^{-14}$ among the three algebraic
routes; the Fourier identity in \eqref{eq:kernel} holds to $10^{-16}$ at
sampled frequencies. The verification suite, run during preparation,
caught and corrected a factor-$2\pi$ normalization error in an
intermediate analytic step---which is its purpose. \pP{}
\end{remark}

\begin{remark}
Equation \eqref{eq:kernel} is consistent with the $\cosh^{-2}$ kernels of
perturbative modular theory in the literature
\cite{SarosiUgajin2018,Lashkari2016,FLPW2016}; it is derived here from
first principles so that no downstream result depends on a cited
normalization. For the field-theoretic applications below,
$\rho_0^{\pm1/2}$ insertions are Euclidean half-modular rotations; the
manipulations are performed at the level of matrix elements between states
of finite modular energy, where the smearing converges absolutely
(exponential decay of $F$ and of matrix elements in $\kappa$); domain
subtleties for unbounded $X$ are handled as standard in this literature.
\pP{} with the stated qualification.
\end{remark}

\section{Spectral Decomposition of the Local/Nonlocal Split}

\begin{definition}[Zero-mode local term]
\label{def:zeromode}
Let $P_0$ denote projection onto modular frequency $\kappa=0$. Define
$\delta K_B^{\mathrm{loc}}\equiv-\,P_0X$ (the $\kappa=0$ component, on
which the multiplier equals $1$) and
$\delta K_B^{\mathrm{nl}}\equiv\delta K_B-\delta K_B^{\mathrm{loc}}$, the
$\kappa\neq0$ part, weighted by $\kappa/(2\sinh(\kappa/2))$.
\end{definition}

\begin{lemma}[Consistency with the null-plane and wedge limits]
On null-plane algebras the vacuum is Markov
\cite{CasiniTesteTorroba2017}; modular Hamiltonians of cuts are exactly
local, all nonzero-frequency cross components vanish, and
Definition~\ref{def:zeromode} reproduces the boost-weighted local term of
\cite{FLPW2016} used in \cite{McGwier2026b}. The definitions agree where
both exist. \pP{}
\end{lemma}

Condition NL is thereby a spectral statement: \emph{it holds iff the
deformation's cross matrix elements with EC-dual perturbations are
supported at modular frequency zero.}

\section{Main Theorem: The Metric Sector is Unconditional}

At order $O(\beta\lambda)$ the cross expectation is bilinear: one
insertion of $J_5$ (from $\delta K_B^{(\beta)}$, smeared by
\eqref{eq:kernel} with the source $\beta$) against one insertion of the
operator creating $|\delta\Psi\rangle$. States dual to linearized EC
solutions are created by boundary stress-tensor modes (gravitons),
axial-current modes (the bulk axial field), and matter primaries.

\begin{theorem}[Unconditional metric sector]
\label{thm:metric}
The graviton channel of Condition NL vanishes identically:
$\langle0|\delta K_B^{\mathrm{nl}}\,|\delta\Psi_T\rangle=0$ for every
stress-tensor perturbation $|\delta\Psi_T\rangle$, every ball, and every
smooth source $\beta$. Consequently the finite-ball torsionful first law
is \emph{unconditionally} equivalent to the linearized Einstein--Cartan
equations in the metric sector, strengthening Theorem 5.2 of
\cite{McGwier2026b}.
\end{theorem}

\begin{proof}
Every term in the cross expectation is a vacuum two-point function of
$J_5$ (spin~1, $\Delta=3$) against $T$ (spin~2, $\Delta=4$), evaluated at
(possibly modular-complexified) points. In any CFT, the vacuum two-point
function of primaries of unequal spin or unequal dimension vanishes; the
modular smearing in \eqref{eq:kernel} acts within the conformal
representation and cannot generate a nonvanishing overlap from a vanishing
one (it multiplies matrix elements by scalars frequency by frequency).
Hence every term is zero. \pP{}
\end{proof}

\begin{remark}
The same selection rule kills the cross channel with any matter primary
not carrying the exact quantum numbers of $J_5$ ($\Delta=3$, spin~1,
axial). The entirety of Condition NL therefore lives in one channel.
\pP{}
\end{remark}

\section{The Surviving Channel: A Universal Kinematic Residue}

\begin{proposition}[Universality]
\label{prop:universal}
In the axial channel, the cross expectation is
\begin{equation}
\mathcal N_B[\beta;\zeta]
=C_J\;\mathcal I_B[\beta;\zeta]\,,
\label{eq:residue}
\end{equation}
where $C_J$ is the two-point normalization of the conserved current,
$\langle J_5^\mu(x)J_5^\nu(0)\rangle
=C_J\,\big(\eta^{\mu\nu}-2x^\mu x^\nu/x^2\big)/x^{6}$ in $d=4$, $\zeta$
labels the axial mode creating $|\delta\Psi\rangle$, and $\mathcal I_B$ is
a purely kinematic functional: the nonzero-modular-frequency part of the
smeared two-point overlap,
\begin{equation}
\mathcal I_B[\beta;\zeta]
=-\!\int_{-\infty}^{\infty}\!\!ds\,\frac{\pi}{2\cosh^2(\pi s)}
\Big[\big\langle
\sigma_s\big(\widetilde{J_5[\beta]}\big)\,J_5[\zeta]
\big\rangle_{\kappa\neq0}\Big],
\label{eq:integral}
\end{equation}
with $\widetilde{\,\cdot\,}$ the half-modular (Euclidean angle $\pi/2$)
conjugation from Proposition~\ref{prop:kernel} and the ball modular flow
given explicitly by the Casini--Huerta--Myers conformal map
\cite{CHM2011}. $\mathcal I_B$ is identical for every CFT possessing a
conserved axial current: dynamics enters \eqref{eq:residue} only through
the single constant $C_J$.
\end{proposition}

\begin{proof}
Conformal symmetry fixes the conserved-current two-point function up to
$C_J$ \pP{}; the smearing, the CHM flow, and the frequency projection are
kinematic operations. \pP{}
\end{proof}

\begin{proposition}[Absorbability: the first law survives either way]
\label{prop:absorb}
If $\mathcal I_B\equiv0$, Condition NL holds and the finite-ball
torsionful first law of \cite{McGwier2026b} holds as stated. If
$\mathcal I_B\neq0$, then---being linear in $\beta$, linear in the mode,
universal, and purely kinematic---it defines a finite, state-independent
bilinear form that can be absorbed as a renormalization of the local term
(equivalently, an improvement of the holographic dictionary relating the
boundary axial source to the bulk contortion boundary value), after which
the improved first law holds. In neither case is the torsionful first law
for finite balls false; the alternative is only between the naive and the
improved dictionary. \pP{} at the stated structural level: the absorption
redefines $\delta K^{\mathrm{loc}}$ by the $\kappa\neq0$ bilinear
evaluated on the source, which is admissible because Condition NL is the
\emph{only} place the split enters the derivation of
\cite{McGwier2026b}, Theorem 5.2.
\end{proposition}

\section{Status, and the One Remaining Number}

\begin{enumerate}[itemsep=2pt]
\item Metric sector: torsionful first law for finite balls
\emph{unconditional} (Theorem \ref{thm:metric}). \pP{}
\item Axial sector: reduced to the single kinematic functional
$\mathcal I_B$ of \eqref{eq:integral}, universal across CFTs; the first
law holds either as stated (if $\mathcal I_B=0$) or in
dictionary-improved form (Proposition \ref{prop:absorb}). \pP{}
\item Evaluation of $\mathcal I_B$: open. \pO{} Two attack paths:
(a) analytic, in $d=2$, where the chiral analogue of the deformation is a
current--current marginal coupling and modular flow of an interval is
exactly geometric on null lines; the null-plane collapse mechanism of
\cite{CasiniTesteTorroba2017} suggests $\mathcal I=0$ there \pC{};
(b) numerical, in $d=4$: \eqref{eq:integral} is a finite-dimensional
quadrature once the CHM flow and the two-point structure are inserted;
the exponential $\cosh^{-2}$ damping and the exponential decay of
$\kappa\neq0$ matrix elements make it numerically benign. This is the
designated computation of the next note.
\item Consequence for the program: with Theorem \ref{thm:metric}, the
gravitational (metric) certification of the New Standard Model
\cite{McGwier2026c} in the presence of torsion sources no longer rests on
any unproved condition; the residual question concerns only the axial
dictionary's finite normalization. \pP{}
\end{enumerate}

\section*{Certification}
Proposition \ref{prop:kernel} is derived in full and machine-verified by
four independent computations (Remark \ref{rem:verification}), with the
verification suite retained alongside the source. Theorem
\ref{thm:metric} rests on the conformal selection rule for two-point
functions, a textbook-grade fact used at \pP{}. Propositions
\ref{prop:universal}--\ref{prop:absorb} are structural and derived at the
level stated. The single remaining unknown, $\mathcal I_B$, is displayed
explicitly and tagged \pO{}; nothing downstream assumes its value.

\section*{Acknowledgments}
The author acknowledges the use of Claude (Anthropic, claude.ai) for
assistance in derivation drafting, numerical verification, and \LaTeX{}
preparation. All technical claims and epistemic classifications were
reviewed and verified by the author, who bears sole responsibility for
the content.

\begin{thebibliography}{9}\small

\bibitem{CHM2011}
H.~Casini, M.~Huerta and R.~C.~Myers, ``Towards a derivation of
holographic entanglement entropy,'' JHEP \textbf{05}, 036 (2011),
arXiv:1102.0440.

\bibitem{FLPW2016}
T.~Faulkner, R.~G.~Leigh, O.~Parrikar and H.~Wang, ``Modular Hamiltonians
for deformed half-spaces and the averaged null energy condition,''
JHEP \textbf{09}, 038 (2016), arXiv:1605.08072.

\bibitem{Lashkari2016}
N.~Lashkari, ``Modular Hamiltonian of excited states in conformal field
theory,'' Phys.\ Rev.\ Lett.\ \textbf{117}, 041601 (2016),
arXiv:1508.03506.

\bibitem{CasiniTesteTorroba2017}
H.~Casini, E.~Test\'e and G.~Torroba, ``Modular Hamiltonians on the null
plane and the Markov property of the vacuum state,''
J.\ Phys.\ A \textbf{50}, 364001 (2017), arXiv:1703.10656.

\bibitem{SarosiUgajin2018}
G.~S\'arosi and T.~Ugajin, ``Modular Hamiltonians of excited states, OPE
blocks and emergent bulk fields,'' JHEP \textbf{01}, 012 (2018),
arXiv:1705.01486.

\bibitem{McGwier2026a}
R.~W.~McGwier, ``Entanglement as the origin of spacetime curvature,''
v4, Zenodo (2026), DOI:
\href{https://doi.org/10.5281/zenodo.21821890}{10.5281/zenodo.21821890}.

\bibitem{McGwier2026b}
R.~W.~McGwier, ``Spin-current sources and ball modular Hamiltonians:
resolving the torsionful first law,'' Zenodo (2026), DOI:
\href{https://doi.org/10.5281/zenodo.21824934}{10.5281/zenodo.21824934}.

\bibitem{McGwier2026c}
R.~W.~McGwier, ``The New Standard Model for unified physics,''
Zenodo (2026), DOI:
\href{https://doi.org/10.5281/zenodo.21829536}{10.5281/zenodo.21829536}.

\bibitem{McGwier2026e}
R.~W.~McGwier, ``The Hehl--Datta coefficient from relative-entropy positivity,''
Zenodo (2026), DOI:
\href{https://doi.org/10.5281/zenodo.21836757}{10.5281/zenodo.21836757}.

\bibitem{McGwier2026f}
R.~W.~McGwier, ``Emergent gravity and a new Lagrangian: a theory of everything we know,''
Zenodo (2026), DOI to be assigned.

\end{thebibliography}

\end{document}
